\section{Discussion and Limitations}
\label{sec:discussion}

\paragraph{Placement and replacement are two axes, and both are real.} The literature tunes the replacement of a skipped step and leaves placement to a threshold or a grid; we tuned placement and left the replacement at plain reuse. Measured in one pipeline at a matched budget the two choices are not equal in size but both matter. Replacing the schedule by a threshold that keeps the same number of calls multiplies the loss by three, so placement dominates. But the replacement is not inert either: holding the schedule and the calls fixed and returning a one-step extrapolation of the velocity instead of the last exact output lowers LPIPS from 0.147 to 0.116 on CogVideoX-2B at seventeen frames, a fifth of the error. The axis was mis-measured in our earlier runs, where the predictor was applied over the whole gap rather than one step, and that version does hurt; the conclusion that reconstruction does not matter was an artefact of that implementation. The right statement is that placement sets which errors are admitted and the replacement reduces what each admitted error costs, so a method should do both.

\paragraph{The prediction must be damped, not scaled.} The size of the correction is not free either. Adding the full one-step difference of the output at every skip gives 0.125 LPIPS; adding half of it gives 0.108, and adding three quarters 0.113, so the optimum is a damped half step and it is flat over a factor of 1.5. Scaling the same difference by the number of steps since the anchor, which is the first-order choice and what feature-space Taylor predictors do, gives 0.442: the trajectory is not locally linear enough to extrapolate that far. The law is a property of the sampler rather than of the prompt, and it explains a published failure in one sentence: a predictor that scales with the gap is extrapolating a curvature it has not measured.

\paragraph{The profile is a property of the sampler.} The order of the first eight steps is identical across prompts and the two profiles correlate at 0.94, which is what makes a fixed schedule defensible at all. The mechanism is a matter of inheritance: an error introduced while the latent is still noise is carried by every subsequent step, because each step conditions on the latent that the previous one produced, while an error introduced near the end is observed by a handful of exact evaluations that re-explain it before the decode. The same mechanism predicts the shape of the tail: since the profile decays, no two decisions near the end deserve equal spacing, which is why the gaps grow rather than stay uniform.

\paragraph{What we do not claim.} The head length is not a universal constant. It is read from a measurement that depends on the step count, the scheduler and the backbone, and we re-measure it once per setting rather than transferring it. The budgets we study lie between ten and twenty-one calls, so the regime below ten, where the head can no longer be held exact, is characterised only by the collapse of every rule we tried. The study covers two backbones, CogVideoX-2B and Wan2.1-1.3B, and two frame counts, and we do not claim a wall-clock result: the budget is counted in transformer calls, which fixes the cost of the network but leaves the decoder and the host out of the comparison. Finally, all baselines are re-run inside one implementation, so the absolute values of Table~\ref{tab:main} are not the numbers those papers report; the ordering within one pipeline is what the table supports.
